% Generated draft: requires mathematical and source review.
\documentclass[11pt]{article}
\usepackage[margin=25mm]{geometry}
\usepackage{fontspec}
\setmainfont{Latin Modern Roman}
\usepackage{amsmath,amssymb,mathtools,unicode-math}
\setmathfont{Latin Modern Math}
\usepackage{enumitem,needspace,xurl,hyperref}
\setlength{\parindent}{0pt}
\setlength{\parskip}{5pt}
\setlength{\emergencystretch}{3em}
\newcommand{\N}{\mathbb{N}}
\newcommand{\Z}{\mathbb{Z}}
\newcommand{\Q}{\mathbb{Q}}
\newcommand{\R}{\mathbb{R}}
\newcommand{\C}{\mathbb{C}}
\newcommand{\F}{\mathbb{F}}
\newcommand{\A}{\mathbb{A}}
\newcommand{\PP}{\mathbb{P}}
\newcommand{\E}{\mathbb{E}}
\newcommand{\eps}{\varepsilon}
\newcommand{\dd}{\,\mathrm{d}}
\begin{document}
\section*{UnsolvedMath ID: 30005962}
\textbf{Source problem code:} catalogue entry
\subsection{Definitions and mathematical statement}
The supplied record is an incomplete preamble: it ends with ``So, it may be reasonable to ask:'' before the actual question. Consequently, no complete conjecture or research question can be stated from this record alone. The assertions below are retained as source-reported background, requiring mathematical review.

The source uses the abbreviation SCY3 for a strict Calabi--Yau threefold, but does not supply its precise definition, ground field, or complete geometric hypotheses. These hypotheses are left unresolved. Write \(\operatorname{Aut}(X)\) for the automorphism group of \(X\), and \(c_2(X)\) for its second Chern class whenever defined under the intended hypotheses. For a projective variety \(Y\), let \(N^1(Y)_{\R}\) denote the real vector space of Cartier divisor classes modulo numerical equivalence. For divisor classes on \(X\), set
\[
c_2(X)^{\perp}=\{\alpha\in N^1(X)_{\R}:\langle c_2(X),\alpha\rangle=0\}.
\]
The pairing here is the intersection pairing intended by the source.

The excerpt defines a \(c_2\)-contraction to be a surjective morphism \(\varphi:X\to B\), where \(X\) is a SCY3 and \(B\) is a normal projective variety with \(\dim B>0\), having connected fibers and satisfying
\[
\langle c_2(X),\varphi^*[H_B]\rangle=0
\]
for an ample Cartier divisor \(H_B\) on \(B\). The source states that, by the pseudo-effectivity of \(c_2(X)\), this is equivalent to
\[
\langle c_2(X),\varphi^*\beta\rangle=0
\qquad\text{for every }\beta\in N^1(B)_{\R}.
\]
It further states that such a morphism is never an isomorphism, because simple connectedness of \(X\) implies that \(c_2(X)\) is not the zero functional on \(N^1(X)\). These are supplied assertions, with their intended hypotheses unresolved.

According to the literal factorization direction in Definition 6 of the excerpt, a \(c_2\)-contraction \(\varphi_0:X\to B_0\) is maximal if, for every \(c_2\)-contraction \(\varphi:X\to B\), there exists a holomorphic map \(g:B\to B_0\) such that
\[
\varphi_0=g\circ\varphi.
\]
This direction is retained as supplied and requires checking against the full definition in the report.

Theorem 5 of the excerpt asserts that each of two rigid SCY3 varieties \(X_k\), \(k\in\{3,7\}\), admits a primitive automorphism \(f_k\in\operatorname{Aut}(X_k)\) of positive entropy. The explicit definitions of \(X_k\) are deferred to reference [10], whose bibliographic entry is absent. The meanings of rigidity, primitivity, and entropy are not specified in the supplied passage. The source cites [11] for \(X_3\) and [9, Lemma 4.6] for \(X_7\), and adds the comparison \(d_2(f)>d_1(f)\) for an automorphism of \(X_7\) discussed there. Neither the definitions of \(d_i\) nor the precise scope of that comparison are supplied. The opening claim that there are only two ``such'' SCY3 examples has an antecedent missing from this record.

Theorem 7 of the excerpt asserts that every SCY3 admitting a \(c_2\)-contraction admits a maximal one, and that a SCY3 admits a birational \(c_2\)-contraction if and only if it is isomorphic to \(X_3\) or \(X_7\). It also asserts that the birational \(c_2\)-contraction is maximal, referring to [10] for proof and details.

For \(K\in\{\Q,\R\}\), retain \(\operatorname{Amp}(X)(K)\) as the notation appearing in the extraction. It is unresolved whether this denotes ample classes or whether a closure or another cone notation was lost during extraction. Proposition 8, attributed in the source to Wilson [14], is reported as
\[
|\operatorname{Aut}(X)|=\infty
\quad\Longrightarrow\quad
\bigl(c_2(X)^{\perp}\cap\operatorname{Amp}(X)(\R)\bigr)\setminus\{0\}\ne\varnothing
\]
for a SCY3 \(X\). The source then gives the speaker's qualified knowledge statement that no example is known to them satisfying both
\[
\bigl(c_2(X)^{\perp}\cap\operatorname{Amp}(X)(\R)\bigr)\setminus\{0\}\ne\varnothing,
\qquad
\bigl(c_2(X)^{\perp}\cap\operatorname{Amp}(X)(\Q)\bigr)\setminus\{0\}=\varnothing.
\]
This is a source-specific knowledge qualification, not a verified nonexistence assertion. The record does not state the question that follows, and these displayed conditions do not determine its intended hypotheses or quantifiers.
\subsection{Short English statement}
This record preserves background on strict Calabi--Yau threefolds, automorphisms, and \(c_2\)-contractions, but stops before the actual question. Its reported real-versus-rational cone conditions require clarification of the cone notation and geometric hypotheses. No complete open-problem statement can be recovered from the supplied excerpt.
\subsection{Sources}
\begin{itemize}
\item Supplied local UnsolvedMath registry record: permanent ID \(30005962\), title ``Strict Calabi--Yau Automorphism Preamble,'' exact source problem code \texttt{catalogue entry}. No proposer is identified. The registry supplies the status label ``open'' and the statement status ``unrecoverable''; these are curation labels, not independently verified mathematical facts.
\item Supplied report citation: \emph{Algebraic Geometry: Wall Crossing and Moduli Spaces, Varieties and Derived Categories} (2025), Oberwolfach Reports, DOI 10.4171/owr/2024/32, pp. 1787--1854. The extraction includes the page marker 1823. Source URL and supplied literature\_sources link labelled ``original OWR report'': \url{https://doi.org/10.4171/owr/2024/32}.
\item Supplied literature assessment dated exactly \texttt{2026-08-22}: the record ends before the question following the established examples and therefore supplies no fully specified claim beyond an unresolved research prompt. This is supplied dated curation derived from the OWR material, not independently verified here; time-sensitive status and cited claims require independent verification.
\item The supplied statement-verification note says that the official report confirms the extraction stops immediately before the actual questions, which are preserved in the following record. This verification claim is part of the supplied registry evidence; neither the report nor the following record has been independently examined here.
\item Internal references preserved from the excerpt: [11] for \(X_3\); [9, Lemma 4.6] for \(X_7\); [10] for the explicit definitions of \(X_k\), their characterization, and proof details; Wilson [14] for Proposition 8. Their full bibliographic entries are not supplied, so these reference numbers cannot be resolved from this record alone.
\end{itemize}
\end{document}
